\section{Introduction}

Programming concurrent and distributed systems is notoriously difficult.\\
Communication-centric programming languages such as Erlang or Elixir, and
frameworks like MPI, are based around \emph{message passing} concurrency, where
processes send and receive messages to coordinate computation, rather than
relying on shared memory.
These actor languages have seen much industrial
interest due to their suitability for distribution and compatibility with
failure recovery idioms such as supervision hierarchies~\cite{Armstrong03}.
However, these languages are also
susceptible to issues like communication mismatches and deadlocks which pose
significant problems at runtime, and are also difficult to detect and fix.

% Behavioural type systems such as \emph{session types}
% \cite{hondaTypesDyadicInteraction1993} encode structured interactions between
% multiple participants, and have been extensively researched and applied to
% real-world programming languages
% \cite{fowlerErlangImplementationMultiparty2016}.
% %
% However, due to their roots in process calculi, session types are particularly
% tailored to languages that make use of \emph{communication channels}.
%
%
\emph{Mailbox types}
\cite{deliguoroMailboxTypesUnordered2018,fowlerSpecialDeliveryProgramming2023}
are a novel behavioural type system with a focus on mailboxes, as used in actor
languages like Erlang and Elixir.
Mailboxes are incoming message queues: multiple processes can \emph{write} to
a mailbox, but only the process that owns a mailbox can read from it.

%   While session types are also
%   applicable in this communication model
%   \cite{mostrousSessionTypingFeatherweight2011a,fowlerErlangImplementationMultiparty2016},
%   mailbox types are specifically designed for them, enabling reasoning about the
%   content of mailboxes.

As an example of mailbox typing, consider the following Erlang implementation of
a \emph{future variable}~\cite{fowlerSpecialDeliveryProgramming2023}, which is a
write-once placeholder variable.

\begin{minipage}{0.45\linewidth}
\begin{lstlisting}[language=erlang]
empty_future() ->
	receive
		{put, X} -> full_future(X)
	end.
full_future(X) ->
	receive
		{get, Pid} -> Pid ! {reply, X},
			full_future(X);
		{put, _} ->
      erlang:error("Multiple writes")
	end.
\end{lstlisting}
\end{minipage}
\hfill
\begin{minipage}{0.5\linewidth}
\begin{lstlisting}[language=erlang, firstnumber=12]
client() ->
  Future = spawn(future, empty_future, []),
  Future ! {put, 5},
  Future ! {get, self()},
  receive
    {reply, Result} ->
      io:fwrite("~w~n", [Result])
  end.
\end{lstlisting}
\end{minipage}
\\

The process \textsf{empty\_future} waits to receive a \textsf{put} message
containing the value to hold.
This is a \emph{selective receive}: other messages may be present in the mailbox
but they will not be retrieved. After the \textsf{put} message has been
received, the process then calls \textsf{full\_future} with the provided value
\textsf{X}.
%
In \textsf{full\_future} the process waits to
receive \textsf{get} messages containing the process ID of the requesting
actor, and responds by sending value \textsf{X} to
the sender.
%
Since a future variable should only be set once, the process throws an error if
more \textsf{put} messages are received. Even in small
examples such as this, protocol violations such as unexpected messages,
forgotten replies and self-deadlocks may occur. All of these problems can be
addressed and statically detected using mailbox typing
\cite{deliguoroMailboxTypesUnordered2018,fowlerSpecialDeliveryProgramming2023}.

A mailbox type consists of a \emph{capability} and a \emph{pattern}.  The
capability indicates whether the type describes messages that must be sent
($\mathsf{!}$) or describe messages to be received from a mailbox ($\mathsf{?}$).
A pattern is a commutative regular expression that describes the contents of a
mailbox. For the above example, we can define the following mailbox types:

\smallmath{
  \mathsf{EmptyFuture} \triangleq \mbtin{(\msg{Put}[\mathsf{Int}] \mbcomp
  \mbstarpayload{\msg{Get}}{\mbtout{\msg{Reply}[\mathsf{Int}]}})}
  \qquad
  \mathsf{FullFuture} \triangleq
      \mbtin
        {
            \mbstarpayload
                {\msg{Get}}
                {\mbtout{\msg{Reply}[\mathsf{Int}]}}
        }
    }

Type $\mathsf{EmptyFuture}$ describes a mailbox that may contain
a single \msg{Put} message with payload of type $\mathsf{Int}$, and a
potentially arbitrary amount of \msg{Get} messages.
Each \msg{Get} message contains a mailbox reference that
expects a \msg{Reply} message to be sent.
%
Type $\mathsf{FullFuture}$ is more restrictive in that it describes a mailbox
only containing an arbitrary amount of \msg{Reply} messages.
By using these types, the type system can statically require that exactly one
\msg{Put} message is sent, and can ensure that every \msg{Get} request receives
a \msg{Reply} message.

Of course, a type system can only guarantee correctness if its own specification is correct.
It is therefore crucial to rigorously prove the properties a type system gives.
Verifying that these proofs are indeed correct can be challenging in itself.
Errors or oversights can remain undetected, resulting in false promises about
type checked programs. For example, a mechanization of multiparty session typing
\cite{hondaMultipartyAsynchronousSession2016} uncovered errors nearly a decade
later~\cite{tiroreMultipartyAsynchronousSession2025}.

In this paper, we report on our mechanization of the semantics of mailbox types,
and of the
\emph{Pat}~\cite{fowlerSpecialDeliveryProgramming2023} programming language that
incorporates mailbox types. 
Although Pat includes extensive pen-and-paper proofs of correctness, the proofs
have not been mechanized.
In particular, we have mechanized numerous properties about mailbox types
themselves, and have mechanized the static semantics of Pat.  During our
mechanization, we discovered several oversights in the original definitions and
proofs which have since been fixed.

\begin{figure}[t]
    \small
  \centering
  \begin{minipage}{0.48\textwidth}
    \scalebox{0.7}{
      \begin{tikzpicture}[
        node distance=0.5cm and 0.5cm,
        every node/.style={draw, rectangle, rounded corners, minimum width=2cm, minimum height=0.8cm, align=center},
        dependency/.style={-Latex, thick}
      ]

      % Nodes
      \node (Message) {Message};
      \node[below=of Message] (MailboxPatterns) {Mailbox Patterns};
      \node[below=of MailboxPatterns] (Types) {Types};
      \node[below =of Types] (Environment) {Environment};
      \node[below =of Environment] (Dblib) {Dblib};
      \node[right=of MailboxPatterns] (Syntax) {Syntax};
      \node[above=of Syntax] (Util) {Util};
      \node[below=of Syntax] (TypingRules) {Typing Rules};
      \node[below=of TypingRules] (Substitution) {Substitution};

      % Direct Dependencies
      \draw[dependency] (Message) -- (MailboxPatterns);
      \draw[dependency] (MailboxPatterns) -- (Types);
      \draw[dependency] (Types) -- (Environment);
      \draw[dependency] (Types) -- (Syntax);
      \draw[dependency] (Dblib) -- (Environment);
      \draw[dependency] (Util) -- (Syntax);
      \draw[dependency] (Environment) -- (TypingRules);
      \draw[dependency] (Syntax) -- (TypingRules);
      \draw[dependency] (TypingRules) -- (Substitution);

      \end{tikzpicture}
    }
  \end{minipage}
  \hspace{-1cm}
  \begin{minipage}{0.48\textwidth}
    \begin{center}
      \begin{tabular}{|l | c | c | c|}
        \hline
        Module & LOC & Definitions & Lemmas\\
        \hline\hline
        Message & 177 & 7 & 9\\
        Mailbox Patterns & 828 & 8 & 44\\
        Types & 695 & 35 & 30\\
        Environment & 3142 & 16 & 131\\
        Util & 48 & 6 & 2\\
        Syntax & 180 & 16 & 0\\
        Typing Rules & 317 & 4 & 12\\
        Substitution & 1325 & 2 & 28\\
        \hline
                & 6712 & 94 & 256\\
        \hline
      \end{tabular}
    \end{center}
  \end{minipage}
  \caption{Dependency graph of the modules in the Rocq project. Transitive dependencies are omitted.}
  \label{fig:rocq-modules}
\end{figure}

We use the Rocq~\cite{teamRocqProver2025} proof assistant for our mechanization.
Figure \ref{fig:rocq-modules} summarizes the structure of our Rocq project and
shows the lines of code, number of definitions and lemmas for each file. 
Due to the necessity of \emph{type combination} and \emph{environment
subtyping}, the \coqe{Environment} module contains, by far, the largest number
of lemmas.

Throughout this paper, definitions and lemmas that are accompanied
with the \rocqlogo -symbol link to the corresponding code in the online
documentation.
The project can be found at \url{https://github.com/edgardSchi/mailbox-types-rocq}.

\subsubsection{Methodology.}
Our mechanization includes both the syntax and type system of Pat, and prove
several properties of the language such as the substitution lemma.
%
We limit the scope of this paper to Pat's static semantics, which alone requires
a significant mechanization effort; a full account of Pat's operational
semantics and preservation theorem would require significant additional work and
is left as a future direction.  We concentrate on properties that are needed for
proving preservation, but also do not require a description of Pat's operational
semantics.

We use Rocq due to its maturity and active development: several successful
mechanizations have been conducted in Rocq, such as \emph{CompCert}
\cite{leroyFormalVerificationRealistic2009}, the first fully formally verified
real-world compiler, and formalizations of session types
\cite{tiroreMultipartyAsynchronousSession2025} and process calculi
\cite{ambalHOPiCoq2021}, modelling communication between processes.

\subsubsection{Contributions.}

The main contributions of this paper are as follows:

\begin{itemize}[leftmargin=*, labelsep=0.5em]
  \item The first mechanization of the syntax and semantics of mailbox types (\S\ref{sec:sup-mailbox-types}).
  \item A formalization of the static semantics of Pat, including definitions that are more
    amenable to mechanization (\S\ref{sec:pat-programming-language}).
  \item Proofs of metatheoretical properties of the static semantics, including
      the nontrivial substitution lemma (\S\ref{sec:pat-programming-language}), requiring over 200 auxiliary lemmas.
  \item We identify missing cases in the original definitions and places where
      the original proofs could be improved, that have since been
      fixed~(\S\ref{sec:sup-mailbox-types} and \S\ref{sec:pat-programming-language}).
\end{itemize}
